\oldpage[IV-2]{40}
\subsection{Descomposiciones irredundantes.}

\begin{proposition}[3.2.1]
\label{IV.2.3.2.1-fr}
Sean $X$ un preesquema localmente noetheriano y $U$ un abierto denso en $X$.
Las condiciones siguientes son equivalentes:
\begin{enumerate}
\item[\emph{a)}] $X$ es reducido.
\item[\emph{b)}] El subpreesquema inducido sobre $U$ es reducido y $X$ carece de
ciclo primo asociado inmerso.
\item[\emph{c)}] $X$ carece de ciclo primo asociado inmerso y, para todo punto
genérico $x$ de una componente irreducible de $X$, se tiene
$\operatorname{long}(\mathcal O_x)=1$.
\end{enumerate}
Los ciclos primos asociados a $X$ coinciden entonces con las componentes
irreducibles de $X$.
\end{proposition}

Está claro que, si $X$ es reducido, lo mismo ocurre con el subpreesquema
inducido sobre $U$. Además, como la existencia de ciclos primos asociados
inmersos es una cuestión local, basta considerar el caso en que
$X=\operatorname{Spec}(A)$ es afín, siendo $A$ noetheriano. Si $A$ es
reducido, se sabe que los ideales primos minimales de $A$ forman una
descomposición primaria reducida de $(0)$ (Bourbaki, \emph{Alg. comm.},
cap.~IV, \S~2, núm.~5, prop.~10) y son los elementos de
$\operatorname{Ass}(A)$; por tanto, $A$ no tiene ideales primos asociados
inmersos, lo que muestra que \emph{a)} implica \emph{b)}. Es inmediato
que \emph{b)} implica \emph{c)}, pues un punto genérico $x$ de una componente
irreducible de $X$ pertenece a $U$, y por tanto $\mathcal O_x$ es un cuerpo.
Finalmente, \emph{c)} implica \emph{a)}: basta, en efecto, observar que, si
$\mathcal N$ es el nilradical de $\mathcal O_X$, que es un ideal coherente,
$\operatorname{Supp}(\mathcal N)$ no puede contener ninguno de los puntos
genéricos de las componentes irreducibles de $X$ por hipótesis; si
$\operatorname{Supp}(\mathcal N)$ no fuera vacío y $x$ fuera uno de los puntos
maximales de este conjunto cerrado, el criterio (3.1.3, \emph{c$'$)}) mostraría
que $x\in\operatorname{Ass}(\mathcal O_X)$, y $\overline{\{x\}}$ sería, por
tanto, un ciclo primo asociado \emph{inmerso} de $X$, en contradicción con la
hipótesis; por tanto, $\mathcal N=0$.
